% Follows DEVIATIONS.md entry by entry (dates and substance copied); each item says what changed, why, and how it changes the reading.
The entries below follow \path{DEVIATIONS.md} of the repository in date order. Each says what changed, why, and how it affects the reading of
the results.
\begin{description}
\item[2026-10-26, physics.] The renderer uses a water-equivalent density with per-component ICRU-44 attenuation, an energy-integrating
 detector and a W-anode spectrum; scatter, noise and blur are applied at training time rather than rendered. \emph{Effect:} the v1 renders are a
 simplified physical model; they concern the v1 study only.
\item[2026-10-26, model size.] SwinV2-Tiny at $512$~px instead of the reference SwinV2-Base at $1024$~px, for compute reasons; the nested
 decoder covers the encoder stages with a plain tail to full resolution. \emph{Effect:} comparisons with the reference model mix the
 regulariser with resolution and capacity; the matched no-twin control isolates the regulariser.
\item[2026-10-26, LIDC sources.] Part of the series came from a Kaggle copy, the rest from LUNA16 (checked voxel-identical to DICOM) or the TCIA
 API; the frame is a body crop with a bounded aspect ratio. \emph{Effect:} none on the results; recorded for reproducibility.
\item[2026-10-26, MIMIC-CXR and Node21.] MIMIC-CXR was dropped (not needed, credentialed access). Node21 was not available as a ready copy;
 it was later fetched from its Zenodo record inside a cloud kernel. \emph{Effect:} none.
\item[2026-10-27, usable CTs.] Short scans needed many padded frame rows; the limit was raised to 50 rows and padded rows are set to air.
 \emph{Effect:} more CTs usable for v1; no locked set involved.
\item[2026-10-27, learning rate.] The pilot passed its gate and a check on synthetic validation CTs kept $2\times10^{-4}$ for every arm; no real
 image was used. \emph{Effect:} none.
\item[2026-10-27, SZCH looked at.] The pre-specified comparison with the reference model was looked at on all 75 SZCH pairs after the v1
 ablation. SZCH became a development set (24 development, 51 held-out images) and its audit development evidence only; the confirmatory audit
 moved to VinDr-RibCXR and JSRT. \emph{Effect:} the main deviation of the project; it is why the second plan exists and why only the 51 held-out
 pairs support fidelity claims.
\item[2026-10-27, secondary metrics.] Band-limited SSIM, the fidelity of the removed signal and a high-frequency energy ratio were added after
 the SSIM gap was seen. \emph{Effect:} they explain the gap; they do not replace the pre-specified metrics.
\item[2026-10-27, Qure projections.] The public Qure.ai projections became the v2 training source after their exact additivity was verified;
 split by patient. \emph{Effect:} the weights inherit CC BY-NC-SA 4.0.
\item[2026-10-29, second plan.] SAP\_v2.md was written after the v2 recipe had been developed on development data; held-out SZCH PSNR and SSIM of
 screening runs had been printed by a screening tool. \emph{Effect:} disclosed; no decision used them, and the tool was changed.
\item[2026-10-29, encoder screen and ablation arms.] Added as exploratory, with an encoder learning rate of $0.25\times$ for vision
 transformers fixed in advance; ConvNeXt-V2-T diverged at the full rate and was rerun at $0.25\times$. \emph{Effect:} the rerun extends a rule
 written for another encoder type.
\item[2026-10-29, accidental start.] A dry-run listing started the held-out uncertainty evaluation by mistake; it was stopped after two minutes
 before writing any output. \emph{Effect:} none; the driver now guards that step.
\item[2026-10-31, BS-LDM.] The released BS-LDM weights were run as an extra baseline outside the analysis plan, with a seeded sampler and our SZCH
 preprocessing. \emph{Effect:} descriptive only; its audit sites were drawn before the masks were pinned and are not paired with the others.
\item[2026-10-31, replacement and tie-break.] Rule R2 first promoted RAD-DINO, then the pre-declared tie-break between the two encoders that met
 R2 chose ConvNeXt-V2-T at two seeds each; the matched no-twin control is the same encoder. With three seeds the gap moved 0.0005 outside the
 tie band; the choice was not reopened. \emph{Effect:} both encoders are reported on the locked sets; the selection used no locked number.
\item[2026-11-01, stored weights.] Seven of 27 evaluated checkpoints carry overflowed (infinite) batch-normalisation variances from mixed-precision
 training in one decoder node. \emph{Effect:} the results describe the stored models as they are; cloud and local outputs agree to 0.0024\,dB.
\item[2026-11-03, cloud evaluation and mask pinning.] The last locked-set outputs were computed on cloud GPUs. Because lung masks differ between
 GPUs, an optional mask cache was added to the audit and rib-benchmark code so that every model is scored on the sites of the final tables; the
 sites were checked to be identical file by file. \emph{Effect:} a post-freeze code change that leaves the computation unchanged when the cache is
 not used.
\item[2026-11-03, synthetic-bank path.] A wrong path in the evaluation driver would have loaded no CT for the synthetic evaluation; fixed before
 it ran. \emph{Effect:} none on results.
\item[2026-11-03, BDCXR-3257 duplicates.] The planned hash rule chained most of the paediatric images into one group; replaced, before any model
 ran, by a correlation rule. The 2\,\% lung-mask rule of the audit is applied to the eroded mask (no image affected). \emph{Effect:} the folds
 are balanced; repeated images of one child cannot be detected without identifiers.
\item[2026-11-03, independent recheck.] Two independent reviews of the statistics and the code found and fixed: an interaction bootstrap that
 dropped repeated draws (intervals about 25\,\% too narrow; one exploratory reading changed), an incomplete implementation of the downstream
 analysis (now with the planned bootstrap intervals), detection metrics that depended on the order of tied scores (up to 0.00055), missing Holm
 correction for the real-nodule contrasts, missing intervals for the reliability results (rerun identical), and the overlap source of the JSRT
 audit (the reference model's map instead of the consensus map). The multiplicity rule was stated explicitly. \emph{Effect:} no
 pre-specified verdict changed; all corrected numbers are those of this thesis.
\item[2026-11-03, v1 H1 record.] A recheck entry had described the v1 H1 test as missing; it had been computed on 27 October 2026 on the test CTs
 with all three lesion-aware arms. The analysis of the 300k finals is a replication. \emph{Effect:} both are reported (Table~\ref{tab:v1_h1}).
\item[3--7 November, v3.] \todo{summarise the v3 entries of DEVIATIONS.md (3 to 7 November) here}
\item[2026-11-08, CT micro-batches.] Gradient accumulation was added to the CT stage after the freeze so that four volumes of 24 slices fit a
 T4 GPU; each step still averages the same four volumes and the network has no batch statistics, so the step is unchanged. \emph{Effect:}
 none on the computation; the frozen hash of the CT code no longer matches the file.
\item[2026-11-08, accounts.] At the author's request one cloud account took the place of another; it prepared its own NIH copy, which was
 compared file by file with the common copy before use. \emph{Effect:} none.
\item[2026-11-09, CT teacher outputs and seed 2.] The design keeps everything derived from CT-RATE in the cloud except metrics and heat-map
 frames, but its own account plan needs the teacher outputs (logits and pooled features of the subset volumes, no images) on other accounts,
 which is possible only through a download; the standard download also brought the teachers' weights and predictions to the author's
 computer. Teacher seed 2 collapsed (constant logits; no code or data fault found). Seeds 5 and 6 were trained with the unchanged recipe
 under a rule fixed before any P1 model of seed 2 (Section~\ref{sec:xm_models}), and seed 5 took slot 2. \emph{Effect:} a disclosed
 conflict with the design's data rule (nothing was released); the CT recipe is seed-unstable (Section~\ref{sec:xm_discussion}).
\item[2026-11-09, last kernels.] The weights of the distilled projection models moved through the author's computer to the only account
 holding CT-RATE, for the cross-modal maps and the VinBigData zero-shot transfer (the same conflict as above). The DenseNet-121 reference runs
 were copied to an account with the common NIH copy, because one account's soft image differed and the rib-emphasis test reads it. The
 pointing-game launcher gained a run filter that leaves the v3 behaviour unchanged. A zero-shot run that failed because it started before its
 model was written was rerun unchanged. VinBigData fine-tuning, planned only if the quota allowed, was not run (about 50 GPU hours).
 \emph{Effect:} none on the hypotheses; one descriptive analysis is missing.
\item[2026-11-10, collection and descriptive summaries.] A rerun that sat in the cloud queue for nine hours was resubmitted unchanged. CT
 teacher seed 5's predictions, first downloaded without them, were fetched like those of the other seeds. The pointing-game summary, the
 grouping of the per-class DeLong tests and the conformal analysis were fixed after the hypothesis results were known
 (Section~\ref{sec:xm_stats}). \emph{Effect:} these descriptive results are post hoc in form; no hypothesis test was affected.
\end{description}
